% =============================================================================
% 6.2 未来展望
% 草稿版本：v0.3 (2026-05-15)
%   v0.1：作为原 5.4.1 \subsubsection*{局限性回顾} + 5.4.2 \subsubsection*{未来工作} 起草。
%   v0.2：按导师意见 3 拆出，分别提级为 6.2 \subsection{研究局限性} 与 6.3 \subsection{研究前景与未来工作}。
%   v0.3：按用户意见 2026-05-15 合并为单一 \subsection{未来展望}，把"局限性 +
%        对应改进"按一一对应关系组成 4 段，删除原过渡话，末段保留多智能体
%        协作开放课题。\label{sec:future-work} 保留以维持其他章节的 \ref 引用。
% 字数：约 1100 字
% 风格规则：v0.2 风格（无 \textbf / \textit / 引例括号 / 双引号 / 破折号）
% 引用：\cite{ref29}（对话状态切换）, \cite{ref1, ref19}（多智能体协作）
% 图表：无
% 依赖宏包：booktabs / tabularx（主稿已加载）
% =============================================================================

\subsection{未来展望}
\label{sec:future-work}

\ref{chap:evaluation} 章的 5 组评测在揭示有效性的同时也暴露了若干局限性，
本节按局限性来源归纳为四类，针对每一类给出对应的工程改进方向，最后给出
超出本文研究范围的多智能体协作扩展作为开放课题。

第一类是多轮长程依赖能力不足。\ref{sec:intent-eval} 节意图识别评测中
\texttt{context\_\bk{}inference} 类 3 条用例端到端通过率仅为 33.3\%，与
单轮场景的 100\% 之间形成 66.7 个百分点差距。当查询同时切换意图与隐式
继承多个槽位时，大语言模型倾向于把所有未在 query 中显式出现的字段都视为
重置而非继承，单纯把上下文摘要拼接到 prompt 不足以实现稳定的多轮槽位
继承。改进方向是在 \texttt{recognizer} 与 \texttt{completer} 之外引入
独立的 \texttt{DialogState} 数据结构，每轮对话结束后把当前
\texttt{locations}、\texttt{time}、\texttt{intent} 三个槽位写入状态对象，
下一轮 prompt 中以结构化 JSON 形式注入，要求大语言模型对每个槽位显式
输出继承或覆盖的二选一决策，把当前的隐式推理转化为显式标注。同时把
多轮评测集从目前的 3 条扩展到 30 条以上，覆盖意图切换、地点切换、
时间切换、复合切换等典型对话状态切换案例 \cite{ref29}。

第二类是数值阈值的精确归属仍依赖确定性分级器兜底。\ref{sec:rag-eval}
节通路 A 检索评测显示，即使在最优混合权重 $\alpha=0.9$ 下，60 条评测集
上仍有约 9 条 Top-1 脱靶，且全部集中在邻近阈值边界、术语别名冲突、跨
类别多文档三类场景，混合检索的工程上限在 85\% 左右，难以单纯依靠权重
调整突破。本文通过通路 B 即确定性 if-else 分级器加 \texttt{grade\_\bk{}id}
硬链接的方式覆盖了这一盲区，但本质上是把问题绕过而非解决。改进方向
是把当前覆盖 GB/T 28592、GB/T 28591、GB/T 27964、GB 50736、GB/T 21987
等 17 项标准共 62 条的知识库扩展到行业气象领域，例如航空气象 ICAO
Annex 3、农业气象作物物候期分级、海洋气象浪高等级等行业标准，对应同步
扩展通路 B 分级器规则与通路 A 评测集。同时引入更细粒度的元数据，例如
\texttt{applicable\_\bk{}region} 区域适用性、\texttt{publish\_\bk{}year}
发布年份等，支持随时间漂移的标准更新。

第三类是受限沙箱的安全性是轻量原型而非生产级方案。\ref{sec:code-eval}
节的沙箱仅做了 builtins 与模块的白名单约束加墙钟超时加结构化错误捕获
四项，能够阻止意外死循环与文件读写，但不能抵御主动恶意攻击。改进方向
是把沙箱升级为进程级隔离，具体有两条候选路径：一是 RestrictedPython
加 PyPy 沙箱，可以在不引入容器开销的前提下实现更严格的字节码级访问
控制；二是 Pyodide WebAssembly 运行时，把代码执行隔离到 WASM 沙箱中，
天然不暴露任何系统调用接口。两条路径在延迟、安全性、可移植性三个维度
上各有取舍，需要进一步实测对比。本文限定在毕业设计研究范围内主动
放弃了这一安全维度。

第四类是模块层与集成层之间的指标衰减。\ref{sec:e2e-eval} 节端到端评测
显示，通路 B 单测 \texttt{citation} 真实性已达 100\%，但端到端集成后
衰减到 26.7\%，意味着大语言模型在最终报告生成阶段仍可能省略检索到的
GB/T 编号。这是一类只有在集成层评测中才会暴露的问题，单独优化任何一个
模块都无法消除。改进方向包含两个方面：一是在 RAG 工具返回结果后立即
把 \texttt{must\_\bk{}cite} 关键字以结构化 JSON 形式注入到模型上下文中，
并在 system prompt 中强制要求最终报告包含该关键字；二是在报告生成完成
后引入独立的引用校验后处理器，对照 \texttt{must\_\bk{}cite} 进行字符串级
匹配，未命中时触发模型二次生成。这两类机制的有效性需要在扩展评测集
上独立验证。

进一步的扩展方向包括把单智能体扩展为多智能体协作架构
\cite{ref1, ref19}，将检索专家、桥接专家、决策专家、报告专家解耦为独立
agent，引入 agent 间共享的 \texttt{Blackboard} 状态空间。多智能体方案
在分工专业化与可调试性上具有优势，但需要解决状态同步与额外延迟两类
工程问题。这一方向超出本文研究范围，保留为后续工作的开放课题。
